\section*{Capítulo \ref{ch:b2:continuous-reactors} --- Reatores contínuos e processos industriais}

\begin{solution}{exo:b2:continuous-reactors:1}
$\tau = V/Q_v = 2000/0.50 = \qty{4.0e3}{s}$, isto é, \qty{1.1}{h}.
\end{solution}

\begin{solution}{exo:b2:continuous-reactors:2}
$k\tau = 1.2 \times 10^{-3} \times 1800 = 2.16$. Tanque: $X = 2.16/3.16 = 0.68$;
tubo: $X = 1 - \mathrm e^{-2.16} = 0.88$.
\end{solution}

\begin{solution}{exo:b2:continuous-reactors:3}
Tubo: $\tau = \ln 100/0.020 = \qty{230}{s}$. Tanque: $X/(1 - X) = 99 = k\tau$,
$\tau = \qty{4950}{s}$, mais de vinte vezes mais longo: em alta \omterm{def:b2:continuous-reactors:conversion}{conversão}, o
tanque agitado paga caro por trabalhar na concentração de saída.
\end{solution}

\begin{solution}{exo:b2:continuous-reactors:4}
Por litro: $120 \times 0.50 = \qty{60}{kJ}$ liberados, aquecendo
\qty{4.2}{kJ/K}: $\Delta T_{\text{ad}} = \qty{14}{K}$. Uma falha do resfriamento
aqueceria a mistura de modo sensível, mas não perigoso, a menos que o solvente
esteja perto de seu ponto de ebulição ou que uma segunda reação comece nessa faixa.
\end{solution}

\begin{solution}{exo:b2:continuous-reactors:5}
Tanque: $Da = X/(1 - X)^2 = 0.8/0.04 = 20$, $\tau = 20/0.010 = \qty{2000}{s}$.
Tubo: $Da = X/(1 - X) = 4$, $\tau = \qty{400}{s}$. Razão 5, contra $4/\ln 5 =
2.5$ para a ordem 1: quanto maior a ordem, mais o tanque agitado sofre com
sua baixa concentração.
\end{solution}

\begin{solution}{exo:b2:continuous-reactors:6}
$k\tau = 5.0$. Três tanques: $X = 1 - (1 + 5/3)^{-3} = 0.947$. Um tanque: $5/6 =
0.833$. Tubo: $1 - \mathrm e^{-5} = 0.993$.
\end{solution}

\begin{solution}{exo:b2:continuous-reactors:7}
$X = 0.70$ significa $k\tau = 0.70/0.30 = 2.33$. Dobrar a vazão divide $\tau$
por dois: $k\tau = 1.17$, $X = 0.54$. Para recuperar 70\,\%, o \omterm{def:b2:continuous-reactors:residence-time}{tempo de
residência}, e portanto o volume, deve ser dobrado.
\end{solution}

\begin{solution}{exo:b2:continuous-reactors:8}
Liberado: $80 \times 2.0 \times 1.0 \times 0.90 = \qty{144}{kW}$. Levado pela
corrente de saída, aquecida de \qty{330}{K} a \qty{350}{K}: $4.0 \times 1.0 \times 20 =
\qty{80}{kW}$. A camisa deve retirar a diferença, \qty{64}{kW}.
\end{solution}

\begin{solution}{exo:b2:continuous-reactors:9}
Dimensões lineares $\times 10$: volume $\times 1000$, área da parede $\times 100$,
área por unidade de volume $\div 10$. O calor liberado cresce com o volume; o
calor retirado através da parede, com sua área: o reator grande retira, por
litro, dez vezes menos calor para a mesma diferença de temperatura.
\end{solution}

\begin{solution}{exo:b2:continuous-reactors:10}
Tubo: $\tau = \ln(0.05/0.10)/(0.05 - 0.10) = \qty{13.9}{min}$, $c_B/c_{A,0} =
\frac{0.10}{-0.05}(\mathrm e^{-1.386} - \mathrm e^{-0.693}) = 0.50$. Tanque: $\tau =
1/\sqrt{0.005} = \qty{14.1}{min}$, $c_B/c_{A,0} = 1.414/(2.414 \times 1.707) =
0.34$. No tanque, A recém-chegado e B já formado estão misturados no mesmo
volume: algum B sempre fica tempo bastante para ser convertido em C enquanto
algum A mal acabou de chegar; no tubo, todas as fatias têm a mesma história.
\end{solution}

\begin{solution}{exo:b2:continuous-reactors:11}
Elevar a temperatura do fluido de resfriamento (e da alimentação) desloca a
reta de retirada para a direita. A interseção fria sobe lentamente até que a
reta se torna tangente à curvatura inferior do S; além disso, o estado frio
desaparece e o tanque salta para o ramo quente (ignição). Abaixando de novo a
temperatura, o tanque permanece no ramo quente até que a reta se torne
tangente à curvatura superior, e só então volta a cair (extinção), a uma
temperatura mais baixa que a da ignição: um ciclo de histerese.
\end{solution}

\begin{solution}{exo:b2:continuous-reactors:12}
Batelada: restam 75\,\% de \qty{2.0}{mol/L}, $1.5 \times 100 = \qty{150}{kJ/L}$,
$\Delta T = 150/4.0 = \qty{37.5}{K}$ acima da temperatura de trabalho, o
bastante para atingir um ponto de ebulição ou uma decomposição. Semibatelada:
no máximo \qty{0.10}{mol/L} sem reagir, $\Delta T = 10/4.0 = \qty{2.5}{K}$. A
adição lenta limita a energia armazenada no reator.
\end{solution}

\begin{solution}{pb:b2:continuous-reactors:1}
\textbf{1.} $r = k[\text{éster}][\ce{OH-}]$; alimentações iguais e uma
estequiometria 1 : 1 mantêm as duas concentrações iguais, $r = kc^2$.
\textbf{2.} $1/c - 1/c_0 = kt$ com $c = 0.10c_0$: $kc_0t = 9$, $t = 9/(0.11
\times 0.10) = \qty{818}{s}$, cerca de \qty{14}{min}.
\textbf{3.} $t_{1/2} = 1/(kc_0) = \qty{91}{s}$; na ordem 2, a velocidade cai
com o quadrado da concentração; por isso, o tempo para reduzi-la à metade
depende do ponto de partida.
\textbf{4.} $Da = kc_0\tau = 0.011\,\tau$ ($\tau$ em segundos).
\textbf{5.} $Q_vc_0 - Q_vc - kc^2V = 0$, isto é, $c_0 - c = k\tau c^2$.
\textbf{6.} Com $c = c_0(1 - X)$: $c_0X = k\tau c_0^2(1 - X)^2$.
\textbf{7.} $Da = 0.90/0.10^2 = 90$.
\textbf{8.} $\tau = 90/0.011 = \qty{8.2e3}{s}$, \qty{2.3}{h}.
\textbf{9.} $V = 0.50 \times 8182 = \qty{4.1e3}{L}$, \qty{4.1}{m^3}.
\textbf{10.} A concentração de saída, $0.010\,\unit{mol/L}$: a velocidade é,
em todo o tanque, cem vezes menor que na entrada.
\textbf{11.} $Da = X/(1 - X) = 9$, $\tau = \qty{818}{s}$, $V = \qty{409}{L}$.
\textbf{12.} Dez vezes menos volume para o tubo.
\textbf{13.} $c_{j-1} - c_j = k\tau_1c_j^2$, sendo $\tau_1$ o \omterm{def:b2:continuous-reactors:residence-time}{tempo de residência}
de um tanque.
\textbf{14.} $\tau_1 = 850/(3 \times 0.50) = \qty{567}{s}$, $k\tau_1 =
\qty{62.4}{L/mol}$. Resolvendo $62.4c_j^2 + c_j - c_{j-1} = 0$: $c_1 = 0.0328$,
$c_2 = 0.0163$, $c_3 = \qty{0.0100}{mol/L}$: 90\,\% de \omterm{def:b2:continuous-reactors:conversion}{conversão}.
\textbf{15.} O tubo, ou a cascata (dez e cinco vezes menores que um único
tanque). Um único tanque ainda pode ganhar por sua temperatura uniforme, pela
facilidade de limpeza e de controle, e porque aqui o volume custa pouco.
\textbf{16.} $\Delta T_{\text{ad}} = 55\,000 \times 100/(997 \times 4180) =
\qty{1.3}{K}$ (\qty{1.2}{K} a 90\,\%).
\textbf{17.} $55 \times 0.50 \times 0.10 \times 0.90 = \qty{2.5}{kW}$.
\textbf{18.} Não: a corrente de saída leva o calor embora com uma elevação de
cerca de um kelvin; nenhum disparo é possível.
\textbf{19.} Vinte vezes mais concentrada: a elevação adiabática passa a
\qty{26}{K} e o calor a \qty{50}{kW}, o que justifica uma serpentina de
resfriamento; sendo a velocidade $kc_0$ vinte vezes maior, o tanque para
90\,\% seria vinte vezes menor (\qty{0.20}{m^3}).
\textbf{20.} $0.50 \times 0.10 \times 0.90 = \qty{0.045}{mol/s}$, $\times 86\,400
\times 82.03 = \qty{319}{kg}$ de etanoato de sódio por dia.
\textbf{21.} A constante de velocidade cresce com a temperatura (Arrhenius);
logo o volume diminui; o custo é o aquecimento da alimentação e, para outras
reações, uma \omterm{def:b2:continuous-reactors:selectivity}{seletividade} menor e mais vapor.
\textbf{22.} O tanque agitado único deve ter $\boldsymbol{V \approx
\qty{4.1}{m^3}}$.
\end{solution}
